Un sustrato $A$ y un catalizador $E$ forman reversiblemente el intermedio $B$: $A+E\mathrel{\mathop{\rightleftharpoons}^{k_1}_{k_2}}B\xrightarrow{k_3}C+E$. Las constantes son positivas, el volumen es constante y se aplica acción de masas. Inicialmente $A(0)=A_0$, $E(0)=E_0$, $B(0)=C(0)=0$, con $A_0\gg E_0>0$. Para la solución explícita use la aproximación $k_1A(t)\simeq k_1A_0=\lambda$ mientras el consumo de sustrato sea pequeño respecto de $A_0$.

\begin{enumerate}[label=(\alph*),ref=\alph*,leftmargin=*]
\item Dé las ecuaciones diferenciales que rigen el sistema.
\item Obtenga la relación entre las cantidades del intermedio $B$ y del catalizador $E$.
\item Determine las cantidades de los cuatro compuestos con el tiempo en el modelo reducido.
\end{enumerate}
